\section{Introduction}
%% ============================================================

High prices alone do not reveal why they persist. When pricing algorithms charge above the competitive benchmark, they may do so because exploration is incomplete, because function approximation introduces upward bias, or because unilateral price cuts trigger costly future retaliation. Only the third mechanism constitutes dynamic collusive enforcement. A credible audit of algorithmic pricing must therefore separate immediate static payoff effects from future continuation penalties rather than relying on converged price levels alone.

Our economic premise is straightforward: holding demand and cost primitives fixed, the public information channel governs whether coordinated prices can be enforced through continuation threats. In repeated games with imperfect public monitoring, the deterrent power of a trigger strategy depends on the gap between detecting a true deviation and falsely punishing cooperation, scaled by the continuation value at risk. We formalize this relationship in a fixed public trigger benchmark. Importantly, this benchmark is a conditional characterization of dynamic incentives; it does not assume that gradient-based neural learners will converge to such a strategy.

Our empirical question is correspondingly focused: do recurrent PPO pricing agents learn policies whose price levels and responses to unilateral price cuts depend on an action-sensitive public signal? To answer this, we compare an informative public reporting channel against an action-independent replacement channel, retraining both firms from scratch under identical budgets. We then freeze the trained policies and inject paired one-period price cuts to decompose total payoff changes into immediate impact effects and discounted continuation effects.

In our experiments, recurrent PPO agents do not sustain information-mediated enforcement. Average prices converge near the one-shot Nash benchmark under both informative and action-independent channels, and removing action information produces no statistically resolved drop in final prices or gain in welfare. Furthermore, counterfactual price cuts lose money immediately at impact rather than through future retaliation, while action-independent monitoring yields an identically zero impulse response by architectural invariance. Our contribution is an integrated framework combining (i) a corrected analytical benchmark for noisy public triggers, (ii) a counterfactual policy audit paired with a proved negative control, and (iii) a complete ten-seed recurrent PPO evaluation.

\subsection{Research Gap: Imperfect Public Monitoring and Recurrent Policy Auditing}

Classical repeated-game theory shows how firms can sustain supra-competitive prices under imperfect public monitoring when rival actions are hidden and only noisy public signals are observed \citep{green1984noncooperative, abreu1990toward, fudenberg1994folk}. In that literature, deviations are deterred when adverse public signals shift play into lower-value continuation regimes, balancing deviation detection against costly false alarms. In algorithmic pricing, \citet{calvano2020artificial} show that tabular Q-learning sustains supra-competitive prices under perfect price observation, and \citet{calvano2021algorithmic} extend tabular Q-learning to imperfect monitoring adapted from \citet{green1984noncooperative}, finding that discrete-price learners conditioned on short signal histories learn finite-duration ``price wars'' after adverse signals. Meanwhile, recent deep reinforcement learning studies such as \citet{deng2025exploring} compare tabular Q-learning, DQN, and PPO under perfect monitoring, focusing primarily on converged price levels.

Two gaps separate these prior studies from our work:
\begin{enumerate}
\item \textbf{Difference from prior imperfect public monitoring studies:} Classical theorems \citep{green1984noncooperative, abreu1990toward, fudenberg1994folk} characterize equilibrium existence for fixed strategies in infinite-horizon games with contemporaneous signals, while \citet{calvano2021algorithmic} study discrete-grid tabular Q-learning and evaluate punishment via observational price-war trajectories. However, observing high prices or post-shock price dips does not by itself prove that future continuation losses outweigh immediate deviation gains. Unlike prior work, we audit continuous-price neural policies under delayed, two-period-averaged Gaussian reports by decomposing counterfactual price-cut payoffs into \emph{immediate impact gain} versus \emph{discounted continuation value change}, paired with an exact action-independent channel replacement (Proposition~\ref{prop:invariance}) as a proved zero-response negative control. We do not claim to refute \citet{calvano2021algorithmic} or generalize across all learners; rather, we test whether observed policy responses reflect dynamic continuation deterrence or immediate static losses.
\item \textbf{Why a windowed GRU with PPO is used:} In our market, firms observe neither rival prices nor contemporaneous one-period signals. Instead, each public report averages two consecutive latent price states, adds Gaussian noise, and arrives after a one-step queue delay. Because of this aggregation, delay, and noise---the signal-extraction problem of imperfect monitoring \citep{green1984noncooperative, calvano2021algorithmic}---a single scalar report cannot distinguish a transient noise shock from a unilateral price cut. Filtering noise requires conditioning on a sequence of past reports, while continuous prices in $[1,\,2.5]$ and continuous reports make tabular discretization over multi-period histories intractable. We therefore pair continuous-action PPO \citep{schulman2017proximal} with a 128-unit GRU encoding a fixed window of 32 public reports (recomputed from a fresh hidden state at each step) as a compact finite-memory sequence encoder. We use GRU+PPO strictly as a controlled testbed for history-dependent continuous pricing policies, not as a claim that it represents commercial pricing software or outperforms tabular, MLP, or no-regret learners (see Appendix~A.5 and~C).
\end{enumerate}

%% ============================================================
